\section{The host-memory tier: a third regime}
\label{sec:tier}

Adding a host-memory KV tier (vLLM CPU offloading; configured sizes
250/375/500 map to effective capacities 0.42/0.63/0.84M tokens per
replica, Section~\ref{sec:calibration}) to the 8-replica fleet does not remove the
collapse failure at supercritical rates. It transforms the failure
into restore-bound congestion, whose persistence is set by tier
capacity exactly as rule~\eqref{eq:ccpustar} of Section~\ref{sec:tier-theory} predicts \evid{R4}\evid{R11}.

\subsection{Restore-bound congestion}
\label{sec:congestion}

Under the standard perturbation at 0.022 sessions/s - the rate at
which the untiered fleet collapses 3/3 - the size-500 tiered fleet
enters a distinct degraded state in 4 of 6 runs (complete counts in
Section~\ref{sec:sizeseries}). Four discriminators separate it from cold collapse,
jointly present in all four congested runs (two at the 300-min
horizon) \evid{R4}:

\begin{itemize}
\item Hit rate stays high: $h$ 0.70--0.85, with tier-hit fraction $\hext$
  0.62--0.79 and restore share of hits 0.66--0.79. Most hits are tier
  restores, not HBM-resident prefixes.
\item The fleet restore rate is pinned at 150--192k tokens/s, which is
  0.81--1.04 of $N \Br$ for the effective per-replica restore rate
  $\Br = 23$k tokens/s (calibration provenance in Section~\ref{sec:calibration};
  per-replica plateau 21.4--24.0k across all four congested runs at
  both fleet scales while demand exceeds it) - the restore-bound
  throughput of equation~\eqref{eq:murestore}.
\item KV occupancy holds at 0.93--0.94, and the tier ingests new writes
  only (measured offload rate 50--78k tokens/s tracks uncached
  prefill plus generation; restored blocks stay tier-resident).
\item The queue diverges slowly and linearly, $\sim$0.8 requests/min over 195
  min, with no transition to the cold state.
\end{itemize}

Against the cold absorbing state at the same operating point
(Section~\ref{sec:anatomy}), the congested state delivers $\sim$2.2$\times$ the throughput and
$\sim$4$\times$ better TTFT \evid{R4}. The failure is not removed: the queue still
diverges, and the state persisted through both 300-min congested
runs.

\begin{figure}[t]
\centering
\includegraphics[width=\columnwidth]{overlay_restore.png}
\caption{Restore-bound congestion at ($N{=}8$, 0.022), size-500 tier:
measured trajectories ($h$, fleet restore rate, waiting) against DES
bands (5 seeds, 300\,min). DES curves are model output (Section~\ref{sec:models});
the four regime discriminators - pinned restore channel, linear
queue divergence, high $h$, no cold collapse - are reproduced in 5 of
5 seeds. Model quantitative gaps ($h$ 0.95 vs measured 0.70--0.85;
restore share 1.0 vs 0.66--0.79) are stated in Section~\ref{sec:validation}.}
\label{fig:restore}
\end{figure}

\subsection{The pinned band is a state outcome, not a channel capacity}
\label{sec:pinnedband}

Reading $\Br$ as a hard per-replica restore ceiling is falsified by
measurement. The two complete-recovery runs sustained 5-minute fleet
restore rates of 235k and 278k tokens/s - 1.28--1.51$\times$ $N \Br$ -
during their escape, then ran their post-cancellation restore phase
at 0.14--0.56 $N \Br$ with waiting 0.3--7 while KV occupancy drained
from 0.83--0.86 to 0.26--0.31 \evid{R4}. Both measured values are reported;
at $n=2$, no distribution is claimed. The 150--192k band of Section~\ref{sec:congestion}
is therefore a congested-state throughput outcome. The raw DMA link
runs at 151.5\,GB/s = 1.19M tokens/s per replica at $\sim$2\% duty in the
congested state, so per-transfer overhead, not the link, binds
there. The serial FCFS restore channel in the DES encodes the
falsified ceiling interpretation and stands as a documented model
deficiency; its concurrent-restore replacement fails the
recovery-branch screen and is not adopted (Section~\ref{sec:failures}) \evid{R9}.

\subsection{Complete recovery at the boundary: one matched pair}
\label{sec:matchedpair}

At base rate 0.020, inside the untiered probability band, the one
measured tiered run fully absorbs the surge: tier-hit fraction
0.42--0.73 for $\sim$2\,h with waiting 1--9, restore share decaying 0.64 to
0, KV occupancy draining to 0.05, and complete recovery by the end
of the 300-min horizon \evid{R5}. The same-day untiered run at the same
operating point collapsed via the delayed mechanism at $\sim$minute 240
\evid{R3}\evid{R5}. This is one matched same-day pair, $n=1$ per side: it
demonstrates that the tier can convert a collapsing boundary point
into a complete recovery in at least one run, and supports no
recovery-probability claim at 0.020.

\subsection{The size series: capacity sets persistence}
\label{sec:sizeseries}

Same-protocol runs at ($N{=}8$, 0.022) across configured tier sizes
\evid{R11}:

\begin{center}
\small
\begin{tabular}{@{}p{0.17\columnwidth}p{0.23\columnwidth}p{0.42\columnwidth}@{}}
\toprule
size (effective) & outcome count & note \\
\midrule
250 (0.42M) & cold collapse 1/1 & enters congestion, then falls through it: tier-hit fraction decays to 0.01 as the working set outgrows the tier \\
375 (0.63M) & complete recovery 2/2 & one run at elevated realized rate (2.02 requests/s) \\
500 (0.84M) & 4 congested, 2 complete recoveries, of 6 & \\
\bottomrule
\end{tabular}
\end{center}

The series separates the two mechanisms of Section~\ref{sec:tier-theory}. Tier
capacity - the varied parameter - sets regime persistence: a tier
too small to hold the recirculating working set loses its external
hits and falls through congestion into the cold state, and the
250/375 contrast brackets the persistence boundary exactly where
rule~\eqref{eq:ccpustar} places it (0.4--0.7M tokens per replica). The
congested-state throughput level is the pinned restore band of
Sections \ref{sec:congestion}--\ref{sec:pinnedband}. Restore bandwidth itself was not varied, so its
effect on persistence is unmeasured. The 375-versus-500 contrast
(2/2 vs 2/6 complete recoveries) suggests a mild size effect but
leaves it unproven at these counts. Realized request rate does not
order outcomes across the series \evid{R11}.

\begin{figure}[t]
\centering
\includegraphics[width=\columnwidth]{overlay_tiercap.png}
\caption{Tier-capacity pair at ($N{=}8$, 0.022): measured size-250 and
size-500 trajectories with DES capacity-sweep bands (model output,
Section~\ref{sec:models}). The DES reproduces the direction and the
cold-persistence boundary of the pair - 250 falling through to cold,
500 staying congested - but produces no complete recoveries at this
rate; the recovery branch is a documented model failure
(Section~\ref{sec:failures}).}
\label{fig:tiercap}
\end{figure}

The executable size series ends at configured size 500: size 650
does not boot on these hosts (the pinned tier plus an uncapped
sidecar file cache exceeds node-allocatable memory), so no statement
is made about larger tiers \evid{R11}.
